\label{sec:conclusion}
表~\ref{tab:summary} 汇总四个问题的答案；每个数字都可在对应小节与结果文件中逐格核对。

\begin{table}[htbp]
\centering
\caption{四个问题的主要结果汇总（详见对应小节；结果文件由附件 3 的模板生成）}
\label{tab:summary}
\small
\begin{tabular}{@{}>{\raggedright\arraybackslash}p{0.155\linewidth}>{\raggedright\arraybackslash}p{0.46\linewidth}>{\raggedright\arraybackslash}p{0.225\linewidth}>{\raggedright\arraybackslash}p{0.10\linewidth}@{}}
\toprule
问题 & 主要答案 & 交付文件 & 详见 \\
\midrule
1（预热平衡，附录 2）
 & \SI{1800}{\second}：中心 $\valQoneCentreTempEnd\,^\circ$C / $\valQoneCentreMoistEnd$，表面 $\valQoneSurfaceTempEnd\,^\circ$C / $\valQoneSurfaceMoistEnd$；失水 $\valQoneMoistLossPct\%$
 & \nolinkurl{result1.xlsx} & 第~\ref{sec:q1} 节 \\
2（全过程，附录 3）
 & \SI{3}{\hour}：中心 $\valQtwoCentreTempThreeH\,^\circ$C / $\valQtwoCentreMoistThreeH$，表面 $\valQtwoSurfaceMoistThreeH$；失水 $\valQtwoMoistLossPct\%$
 & \nolinkurl{result2.xlsx}、\texttt{result2\_全过程.xlsx} & 第~\ref{sec:q2} 节 \\
3（烘干时长）
 & $t_{\mathrm{end}}=\valQthreeDryHours$ h（$\valQthreeDryDays$ 天；外推 $\valQthreeDryHoursExtrapolated$ h），由\textbf{中心}控制
 & \nolinkurl{result3.xlsx} & 第~\ref{sec:q3} 节 \\
4（尺寸变化）
 & $t_{\mathrm{end}}=\valQfourDryHours$ h，末半径 $\valQfourRadiusAtDryCm$ \si{\centi\meter}；同物性固定半径需 $\valQfourDryHoursFixed$ h
 & \nolinkurl{result4.xlsx} & 第~\ref{sec:q4} 节 \\
\bottomrule
\end{tabular}
\end{table}

\noindent 三条跨问题的结论：

\begin{enumerate}[leftmargin=*,itemsep=0.15em,topsep=0.2em]
  \item \textbf{过程由内部扩散控制。} 传质 Biot 数由 $\valQtwoBiotMassInit$ 升到 $10^3$ 量级，烘干时间对 $h$ 的弹性只有 $\valElasH$、对 $h_m$ 为 $\valElasHm$，而对扩散前因子 $D_0$ 为 $\valElasDzero$、对烘房温度平台为 $\valElasTplateau$。工艺上应提温或减小尺寸，加大风速无效。
  \item \textbf{“各处低于 0.15”是最敏感的一处表述。} 满足它需要 $\valQthreeDryHours$ h，而此时体积平均含水率已低至 $\valQthreeMeanMoistAtDry$ \si{\kilo\gram\per\kilo\gram}；改用平均含水率判据只需 $\valAltQthreeMeanCriterionHours$ h（$\valAltQthreeMeanCriterionPct\%$）。判据的定义比任何参数的 $\pm20\%$ 扰动都更决定答案。
  \item \textbf{数值上的答案是可核查的，物理上的答案还有已知的偏差方向。} 离散与时间积分误差合计远小于 $0.1\%$（第~\ref{sec:validation} 节），解释性选择中最大的两项为守恒写法（$\valAltMassConservativePct\%$）与判据表述（$\valAltQthreeMeanCriterionPct\%$）；题给模型省略了蒸发潜热，界定分析表明真实烘干时间很可能比主结果长 $\valExtCappedPct\%$--$\valExtCapLocalPct\%$（第~\ref{subsec:extension} 小节）。
\end{enumerate}
